\BourbakiExercisesSection

\begin{enumerate}
\def\labelenumi{\arabic{enumi}.}
\tightlist
\item
  设 \(\alpha \mapsto f_{\alpha}\) 是 \(\Omega\) 在 E 上的一个作用。设 F 为 \(\Omega\) 在 \(\mathbf{E}^{\mathbb{E}}\) 中的像，G 为 \(\mathbf{E}^{\mathbb{E}}\) 在律 \((f,g) \mapsto f \circ g\) 下由 F 与 E 的恒等映射生成的稳定子集。
\end{enumerate}

\begin{enumerate}
\def\labelenumi{(\alph{enumi})}
\item
  证明 E 的每个在 \(\Omega\) 的作用下稳定的子集，在 \(E^{E}\) 在 E 上的典范作用限制到 G 上之下也是稳定的（第 1 号，例 3）。
\item
  设 X 是 E 的一个子集；证明由 X 生成的 E 的稳定子集是诸 \(f(x)\) （f 遍历 G，x 遍历 X）所成的集合。
\end{enumerate}

\begin{enumerate}
\def\labelenumi{\arabic{enumi}.}
\setcounter{enumi}{1}
\item
  设 \(\bot\) 是 E 上的一个关于结合的内律 \(\top\) 双重分配的内律；证明，若 \(x \perp x'\) 与 \(y \perp y'\) 在律 \(\top\) 下是可消去的，则 \(x \perp y'\) 与 \(y \perp x'\) 在律 \(\top\) 下可交换（以两种不同的方式计算复合 \((x \top y) \perp (x' \top y')\) ）。特别地，若律 \(\bot\) 有单位元，则律 \(\top\) 下两个可消去的元素在该律下可交换；若 E 的所有元素在律 \(\top\) 下都可消去，则该律是交换的。
\item
  设 \(\top\) 与 \(\bot\) 是集合 \(E\) 上的两个内律，使得 \(\bot\) 关于 \(\top\) 是右分配的。
\end{enumerate}

\begin{enumerate}
\def\labelenumi{(\alph{enumi})}
\item
  若律 \(\top\) 有单位元 \(e\) ，则 \(x \perp e\) 在律 \(\top\) 下是幂等的，对所有 \(x \in \mathbb{E}\) ；若进一步存在 \(y\) 使得 \(x \perp y\) 在律 \(\top\) 下可消去，则 \(x \perp e = e\) 。
\item
  若律 \(\perp\) 有单位元 u，且存在 \(z \in E\) 对律 \(\perp\) 与 \(\top\) 都可消去，则 u 在律 \(\top\) 下可消去。
\end{enumerate}

¶ 4. 设 \(\top\) 与 \(\bot\) 是 E 上的两个内律，各有一个单位元；若由每个律导出的 E 在自身上的左作用关于另一个律是分配的，则 E 的每个元素在这两个律下都是幂等的（若 \(e\) 是 \(\top\) 的单位元， \(u\) 是 \(\bot\) 的单位元，先证明 \(e \perp e = e\) ，注意 \(e = e \perp (u \top e)\) ）。

\begin{enumerate}
\def\labelenumi{\arabic{enumi}.}
\setcounter{enumi}{4}
\tightlist
\item
  在集合 E 上假设给定了三个内合成律：一个加法（不必结合也不必交换）、一个乘法（不必结合）以及一个记作 \(\top\) 的律。假设乘法有单位元 \(e\) ，律 \(\top\) 关于乘法是右分配的，且律 \(\top\) 关于加法是左分配的。证明：若存在 \(x, y, z\) 使得 \(x \top z, y \top z\) 与 \((x + y) \top z\) 在乘法下可消去，则 \(e + e = e\) （利用习题3(a)）。
\end{enumerate}

¶ 6. E 上的一个内合成律 \(\top\) 称为在 E 上确定一个拟群结构，如果对所有 \(x \in E\) ，左平移与右平移 \(\gamma_x\) 与 \(\delta_x\) 都是 E 到其自身的双射映射。称一个拟群是分配的，如果律 \(\top\) 关于其自身是双重分配的。

\begin{enumerate}
\def\labelenumi{(\alph{enumi})}
\item
  确定一个具有 \(n\) 个元素的集合上的所有分配拟群结构，对 \(2 \leqslant n \leqslant 6\) 。
\item
  *证明有理数集 \(\mathbf{Q}\) 配备律 \((x,y)\mapsto\frac{1}{2}(x+y)\) 是一个分配拟群。*
\item
  在分配拟群 E 中，每个元素都是幂等的。由此推出，若 E 含有多于一个元素，则律 \(\top\) 既不能有单位元，也不能是结合的。
\item
  分配拟群 E 的右平移和左平移都是 E 的自同构。
\item
  若E是有限分配拟群，则E的每个稳定子集上诱导的结构都是分配拟群结构。
\item
  设 E 是一个分配拟群。若 R 是一个与律 T 左（相应地，右）相容（第 3 号）的等价关系，则模 R 的类都是 E 的稳定子集。若 E 是有限的，则所有这些类都可以通过左（相应地，右）平移从一个得到另一个；在同一条件下，若 R 与律 T 相容，则 E/R 上的商结构是一个分配拟群结构。
\item
  设 E 是一个分配拟群。E 中与给定元素 a 可交换的元素所成的集合 \(A_{a}\) 是稳定的；若 E 是有限的，则对所有 \(x \in A_{a}\) 有 \(A_{x} = A_{a}\) （注意，利用 (e)，存在 \(y \in A_{a}\) 使得 \(x = a \top y\) ），且当 x 遍历 E 时，诸集合 \(A_{x}\) 与关于一个和律 T 相容的关系的等价类相同。
\item
  若 E 是有限分配拟群且律 \(\top\) 是交换的，则 E 的元素个数是奇数（考察 E 的元素的有序对 \((x,y)\) ，满足 \(x\top y=y\top x=a\) ，其中 a 是给定的）。
\end{enumerate}

\begin{enumerate}
\def\labelenumi{\arabic{enumi}.}
\setcounter{enumi}{6}
\tightlist
\item
  假设在集合 E 上给定一个（结合且交换的）加法，在该加法下 E 的所有元素都可逆，以及一个（结合的）乘法，它关于加法 \(*\) 是双重分配的（换言之，一个环结构） \(_{*}\) ；记 \([x,y]=xy-yx\) ；律 \((x,y)\mapsto[x,y]\) 关于加法是双重分配的。要使 x 与 y 在乘法下可交换，必要且充分的是 \([x,y]=0\) ；我们有恒等式
\end{enumerate}

\[
[ x, y ] = - [ y, x ]; \quad [ x, [ y, z ] ] + [ y, [ z, x ] ] + [ z, [ x, y ] ] = 0
\]

（第二个恒等式称为“雅可比恒等式”）。第二个恒等式也可以写成

\[
[ x, [ y, z ] ] - [ [ x, y ], z ] = [ [ z, x ], y ]
\]

它表达了律 \([x, y]\) 的“结合性偏差”。

\begin{enumerate}
\def\labelenumi{\arabic{enumi}.}
\setcounter{enumi}{7}
\tightlist
\item
  在与习题7相同的假设下，设 \(x \top y = xy + yx\) ；则律 \(\top\) 是交换的且关于加法是双重分配的，但一般不是结合的。
\end{enumerate}

\begin{enumerate}
\def\labelenumi{(\alph{enumi})}
\item
  对所有 \(x \in \mathbb{E}\) ，证明 \(\top^{m+n} x = \left( \begin{array}{c} m \\ \top \end{array} x \right) \top \left( \begin{array}{c} n \\ \top \end{array} x \right)\) 。
\item
  若记 \([x, y, z] = (x \top y) \top z - x \top (y \top z)\) （律 \(\top\) 的结合性偏差），证明恒等式
\end{enumerate}

\[
[ x, y, z ] + [ z, y, x ] = 0
\]

\[
[ x, y, z ] + [ y, z, x ] + [ z, x, y ] = 0
\]

\[
[ x \top y, u, z ] = [ u, [ (x \top y), z) ] ]
\]

（记号 \([x, y]\) 含义与习题7中相同）

\[
[ x \top y, u, z ] + [ y \top z, u, x ] + [ z \top x, u, y ] = 0.
\]

¶ 9. 假设在 E 上给定一个（结合且交换的）加法，使得 E 的所有元素在该加法下都可逆，并给定一个乘法，该乘法非结合，但交换且对加法双重分配。进一步假设 \(n \in \mathbf{Z}\) , \(n \neq 0\) 和 \(nx = 0\) 蕴含 \(x = 0\) 在 E 中成立。证明：若记 \([x, y, z] = (xy)z - x(yz)\) ，则恒等式

\[
[ x y, u, z ] + [ y z, u, z ] + [ z x, u, y ] = 0
\]

成立，则 \(x^{m + n} = x^m x^n\) 对所有 \(x\) 成立（对 \(p\) 用归纳法证明：恒等式 \([x^q, y, x^{p - q}] = 0\) 对 \(1 \leqslant q < p\) 成立）。

\begin{enumerate}
\def\labelenumi{\arabic{enumi}.}
\setcounter{enumi}{9}
\tightlist
\item
  设 E 是一个交换幺半群，其律以加法记法书写，其单位元记为 0。设 \(\top\) 是 E 上的一个关于加法分配的内律，且满足 \(0 \top x = x \top 0 = 0\) （对所有 \(x \in \mathbb{E}\) ）。设 S 是 E 的一个在加法下及在由 \(\top\) 导出的每个外律下都稳定的子集；令 \(\overline{\mathbb{E}}\) 表示 E 关于 S 的加法差幺半群， \(\varepsilon\) 为 E 到 \(\overline{\mathbb{E}}\) 中的典范同态。则在 \(\overline{\mathbb{E}}\) 上存在一种且仅一种律 \(\overline{\top}\) ，它关于加法是分配的，且使得 \(\varepsilon\) 对律 \(\top\) 与 \(\overline{\top}\) 是同态。若律 \(\top\) 是结合的（相应地，交换的），则律 \(\overline{\top}\) 亦然。
\end{enumerate}

